% p1-08-empirical-results

\section{P1 §8. Representative Empirical Results}

8. Representative Empirical Results

   Figure (fig-p1-ch8-empirics). Empirics triptych --- best-fit residual scatter / coverage(C) S-curve /
                                    alpha-phi on the number-line.

  8.1 Best-Fit Representation Functional
  For a target constant X $\in$ T, the best-fit symbolic representative at complexity C is

  F^*(X) := argminF $\in$ S(C) (|F - X|)/(|X|). 27

  This defines a nonlinear projection map PiS(C) : T $\rightarrow$ S(C). The resulting residuals
  {dC(Xi)} constitute the primary empirical observations.

  8.2 Selected Representative Fits
  The following expressions represent low-residual best fits identified by exhaustive
  search over S(C) at complexity bounds C $\leq$ 10. They are reported as empirical results of
  the constrained search, not as theoretical derivations. Several appear in the PhD
  monograph Chapter 34 exploratory catalogue (PhD: Ch.34 Catalogue); their presence
  here is coincidental with that catalogue and does not constitute re-use of the catalogue
  as evidence.

  Strong coupling constant (\citealp{pdg2024}; alphas(mZ) = 0.1180 ± 0.0009):

  alpha$\varphi$ := ($\varphi$-3)/(2) = ((2$\varphi$-3))/(2) approx 0.11803..., dC(alphas) $\sim$ O(10-4), C = 3. 28

  Remark on alpha$\varphi$. The algebraic identity $\varphi$-1 = $\varphi$ - 1 allows rewriting alpha$\varphi$ =
  (1)/(2)($\varphi$-1)3 = (2$\varphi$-3)/2 = (5-2)/2, confirming this is a degree-3 polynomial in a
  quadratic irrational with small integer coefficients --- the lowest-complexity non-trivial
  element of the grammar near alphas(mZ). The relative residual d3(alphas) approx
  (0.11803398875 - 0.1180)/0.1180 approx 2.9 $\times$ 10-4 lies within the experimental
  uncertainty sigma/alphas approx 0.76\%. After Bonferroni correction over |S(3)| =
  1{,}161 trials, the corrected p-value at tolerance 2.9 $\times$ 10-4 must be assessed against
  the null distribution. The specific definition alpha$\varphi$ := $\varphi$-3/2 is used here; the PhD
  monograph Chapter 38 contains an inconsistency between two candidate definitions

  (PhD: Ch.38 AlphaPhi), and only this definition is adopted.

  Inverse fine-structure constant (\citealp{codata2022}; alpha-1 = 137.035 999 177(21)):

  pi4 $\cdot$ $\varphi$4 $\cdot$ e2 / 36 approx 137.036581054, dC(alpha-1) approx 4.25 $\times$ 10-6, C = 10. 29

  The residual approx 4.25 $\times$ 10-6 at this complexity is non-trivial but must survive
  correction over |S(10)| = 75{,}249 trials to be statistically significant. The expression is
  reported as a compact fit, not as an explanation of QED.

  Electroweak mixing angle (\citealp{pdg2024}; sin2theta(MZ) = 0.23129(4)):

  3-1 $\cdot$ pi-2 $\cdot$ $\varphi$4 approx 0.231489, dC(sin2thetaW) approx 8.6 $\times$ 10-4, C = 7. 30

  Solar neutrino mixing (\citealp{pdg2024}; sin2theta12 = 0.303+0.012-0.011):

  8 $\cdot$ $\varphi$-5 $\cdot$ pi $\cdot$ e-2 approx 0.306699, dC(sin2theta12) approx 1.2 $\times$ 10-2, C = 8. 31

        Remark. The reported approximations depend explicitly on the choice of symbolic
        grammar S(C), the complexity cutoff C, and the error functional. They should be
        interpreted as properties of the representation system, not as intrinsic identities
        of physical constants. At C = 8, |S(8)| = 32{,}841; whether the observed residuals
        are meaningfully smaller than a grammar-aware null requires the permutation
        calibration of Section 6.

  8.3 Full 42-Entry Trinity Catalogue Table
  The complete legacy catalogue table is restored here because it is the object referred to
  by     the      name        Catalog42.     The      table     is     imported       from
  t27/research/trinity-pellis-paper/G2TRINITYV1.0FRAGRANCE.tex and cross-referenced
  against t27/proofs/trinity/Catalog42.v. Its evidentiary status is deliberately split into
  two layers:

  - Catalogue layer: 42 rows across 9 physics sectors, each giving a proposed Trinity
  monomial or near-monomial parametrization.
  - Coq layer: Catalog42.v is a Coq aggregation file with source-level Qed markers and no
  Admitted markers in that file; it aggregates a representative verified theorem set, but it
  does not, in its current form, provide one independent theorem named for every one of
  the 42 table rows.

  Therefore the table below is part of the paper's data/provenance layer, while statistical
  significance still depends on the null-model and MDL procedures of Sections 5-7.

  For print readability, the catalogue is printed below as a row-wise audit ledger. Each
  line has the same six fields as the source table: ID, sector, constant, reference value,
  Trinity expression, and residual/tier.

  - G01 --- Sector: Gauge / running; constant: alpha-1 fine structure; reference value:
  137.036; expression: 4$\cdot$9$\cdot$pi-1$\varphi$ e2; residual/tier: 0.029\% V.
  - G02 --- Sector: Gauge / running; constant: alphas(mZ)=alpha$\varphi$; reference value:
  0.11800; expression: $\varphi$-3/2; residual/tier: 0.029\% V.
  - G03 --- Sector: Gauge / running; constant: sin2thetaW; reference value: 0.23121;
  expression: 3-2pi2$\varphi$3 e-3; residual/tier: 0.086\% V.

  - G04 --- Sector: Gauge / running; constant: cos2thetaW; reference value: 0.76879;
  expression: 2pi$\varphi$-2e-1; residual/tier: 0.175\% C.
  - G05 --- Sector: Gauge / running; constant: alphas/alpha2 ratio; reference value:
  3.7387; expression: 2pi$\varphi$ e-1; residual/tier: 0.034\% V.
  - G06 --- Sector: Gauge / running; constant: alpha(mZ)/alpha(0) running; reference
  value: 1.0631; expression: 3$\varphi$2 e-2; residual/tier: 0.017\% V.
  - H01 --- Sector: Electroweak; constant: mH [GeV]; reference value: 125.20; expression:
  4$\varphi$3e2; residual/tier: 0.032\% V.
  - H02 --- Sector: Electroweak; constant: mW [GeV]; reference value: 80.369; expression:
  4$\cdot$3-1pi3$\varphi$-1e; residual/tier: 0.051\% V.
  - H03 --- Sector: Electroweak; constant: mZ [GeV]; reference value: 91.188; expression:
  7$\cdot$3pi-1$\varphi$3e-2; residual/tier: 0.068\% V.
  - H04 --- Sector: Electroweak; constant: GammaZ [GeV]; reference value: 2.4955;
  expression: 4$\cdot$3-1pi$\varphi$ e-1; residual/tier: 0.087\% V.
  - H05 --- Sector: Electroweak; constant: mt/mH ratio; reference value: 1.3784;
  expression: 7pi-1$\varphi$-1; residual/tier: 0.092\% V.
  - H06 --- Sector: Electroweak; constant: mt/mW ratio; reference value: 2.1472;
  expression: 7pi-1$\varphi$2e-1; residual/tier: 0.057\% V.
  - H07 --- Sector: Electroweak; constant: sigmahad at Z [nb]; reference value: 41.48;
  expression: 3pi$\varphi$ e; residual/tier: 0.066\% V.
  - L01 --- Sector: Leptons / Koide; constant: me [MeV]; reference value: 0.51100;
  expression: 2pi-2$\varphi$4e-1; residual/tier: 0.017\% V.
  - L02 --- Sector: Leptons / Koide; constant: mmu [MeV]; reference value: 105.658;
  expression: 8$\cdot$9$\cdot$pi-4$\varphi$2e4; residual/tier: 0.043\% V.
  - L03 --- Sector: Leptons / Koide; constant: mtau [MeV]; reference value: 1776.86;
  expression: 5$\cdot$33pi-3$\varphi$5e; residual/tier: 0.067\% V.
  - L04 --- Sector: Leptons / Koide; constant: ymu/ytau ratio; reference value: 0.05946;
  expression: 3-2pi-1$\varphi$-1e; residual/tier: 0.077\% V.
  - K01 --- Sector: Leptons / Koide; constant: Q(e,mu,tau) Koide; reference value: 0.66667;
  expression: 8$\varphi$-1e-2; residual/tier: 0.370\% C.
  - K02 --- Sector: Leptons / Koide; constant: Q(u,d,s) Koide; reference value: 0.5620;
  expression: 4$\varphi$-2e-1; residual/tier: 0.012\% V.
  - K03 --- Sector: Leptons / Koide; constant: Q(c,b,t) Koide; reference value: 0.6690;
  expression: 8$\varphi$-1e-2; residual/tier: 0.020\% V.
  - Q01 --- Sector: Quark masses; constant: mu [MeV]; reference value: 2.160; expression:
  pi2$\varphi$ e-2; residual/tier: 0.056\% V.
  - Q02 --- Sector: Quark masses; constant: md [MeV]; reference value: 4.670; expression:
  3$\varphi$3e-1; residual/tier: 0.109\% C.
  - Q03 --- Sector: Quark masses; constant: ms [MeV]; reference value: 93.40; expression:
  7pi$\varphi$3; residual/tier: 0.261\% C.
  - Q04 --- Sector: Quark masses; constant: mc [GeV]; reference value: 1.273; expression:
  pi2$\varphi$-4e2; residual/tier: 0.083\% V.

  - Q05 --- Sector: Quark masses; constant: mb [GeV]; reference value: 4.183; expression:
  5pi$\varphi$-2e-1; residual/tier: 0.054\% V.
  - Q06 --- Sector: Quark masses; constant: mt [GeV]; reference value: 172.57; expression:
  4$\cdot$9$\cdot$pi-1$\varphi$4e2; residual/tier: 0.043\% V.
  - Q07 --- Sector: Quark masses; constant: ms/md ratio; reference value: 20.000;
  expression: 8$\cdot$3$\cdot$pi-1$\varphi$2; residual/tier: 0.002\% SG.
  - Q08 --- Sector: Quark masses; constant: md/mu ratio; reference value: 2.162;
  expression: pi2$\varphi$ e-2; residual/tier: 0.038\% V.
  - C01 --- Sector: CKM; constant: |Vus|; reference value: 0.22431; expression:
  2$\cdot$3-2pi-3$\varphi$3e2; residual/tier: 0.051\% V.
  - C02 --- Sector: CKM; constant: |Vcb|; reference value: 0.04100; expression: pi3$\varphi$-3e-1;
  residual/tier: 0.073\% V.
  - C03 --- Sector: CKM; constant: |Vub|; reference value: 0.00394; expression:
  3-2pi-3$\varphi$2e-1; residual/tier: 0.068\% V.
  - C04 --- Sector: CKM; constant: deltaCPCKM [deg]; reference value: 65.9; expression:
  2$\cdot$3$\varphi$ e3; residual/tier: 0.061\% V.
  - N01 --- Sector: PMNS; constant: sin2theta12; reference value: 0.30700; expression:
  8$\varphi$-5pi e-2; residual/tier: 0.089\% V.
  - N02 --- Sector: PMNS; constant: sin2theta23; reference value: 0.546; expression:
  4$\cdot$3-1pi$\varphi$2e-3; residual/tier: 0.085\% V.
  - N03 --- Sector: PMNS; constant: sin2theta13; reference value: 0.02224; expression:
  3pi$\varphi$-3$\cdot$10-2; residual/tier: 0.040\% V.
  - N04 --- Sector: PMNS; constant: deltaCPPMNS [deg]; reference value: 129.1;
  expression: 8pi3/(9e2); residual/tier: 0.037\% V; under formula-audit in Coq comments.
  - M01 --- Sector: Cosmology; constant: Omegab; reference value: 0.04897; expression:
  4$\varphi$-2pi-3; residual/tier: 0.041\% V.
  - M02 --- Sector: Cosmology; constant: OmegaDM; reference value: 0.2607; expression:
  7$\cdot$3-1pi-2$\varphi$3; residual/tier: 0.071\% V.
  - M03 --- Sector: Cosmology; constant: OmegaLambda; reference value: 0.6841;
  expression: 5pi-2$\varphi$2e-1; residual/tier: 0.086\% V.
  - M04 --- Sector: Cosmology; constant: ns spectral index; reference value: 0.9649;
  expression: 3$\varphi$3pi-4e2; residual/tier: 0.094\% V.
  - D01 --- Sector: QCD hadrons; constant: fK [MeV]; reference value: 157.55; expression:
  pi4$\varphi$; residual/tier: 0.039\% V.
  - P01 --- Sector: Loop quantum gravity; constant: gammaBI Barbero-Immirzi; reference
  value: 0.23753; expression: $\varphi$-3=5-2; residual/tier: 0.62\% C.

  8.4 Coq Scope Audit: Why the Number Is Larger Than 42
  The number 42 is the size of the restored legacy formula catalogue, not the size of the
  Coq proof base. A source audit of the available t27 checkout shows a larger formal layer
  around this catalogue:

  - t27/proofs/trinity/: 13 Coq files, 281 theorem-like declarations, 122 source-level Qed.
  markers, 32 source-level Admitted. markers, and 0 Abort. markers.
  - t27/coq/: 10 Coq kernel files, 73 theorem-like declarations, 35 source-level Qed.
  markers, 0 Admitted. markers, and 11 source-level Abort. markers.
  - t27/proofs/: 18 Coq proof files in total, 296 theorem-like declarations, 127 source-level
  Qed. markers, 32 source-level Admitted. markers, and 0 Abort. markers.

  Within t27/proofs/trinity, the following files constitute the topic-specific proof layer for
  the physical-constant programme:

  - CorePhi.v: 13 declarations, 12 Qed., 0 Admitted.; root identities for $\varphi$, including
  $\varphi$2+$\varphi$-2=3.
  - AlphaPhi.v: 13 declarations, 12 Qed., 0 Admitted.; closed form and numeric windows
  for alpha$\varphi$=$\varphi$-3/2.
  - BoundsGauge.v: 21 declarations, 9 Qed., 0 Admitted.; gauge-sector bounds for
  G-series formulas.
  - BoundsMixing.v: 22 declarations, 10 Qed., 0 Admitted.; CKM/PMNS mixing bounds.
  - BoundsMasses.v: 27 declarations, 11 Qed., 0 Admitted.; mass-ratio and Higgs-sector
  bounds.
  - BoundsQuarkMasses.v: 16 declarations, 4 Qed., 4 Admitted.; additional quark-mass
  formula checks.
  - BoundsLeptonMasses.v: 16 declarations, 0 Qed., 8 Admitted.; lepton-mass formula
  checks still requiring closure.
  - FormulaEval.v: 30 declarations, 17 Qed., 0 Admitted.; generic monomial evaluator,
  with a header note stating that it captures the 69-formula Trinity v0.9 framework at the
  evaluation-interface level.
  - DerivationLevels.v: 29 declarations, 17 Qed., 0 Admitted.; L1-L7 derivation-level
  classification and representative derivation paths.
  - ExactIdentities.v: 27 declarations, 5 Qed., 11 Admitted.; Lucas/Pell/Fibonacci identity
  layer.
  - ConsistencyChecks.v: 22 declarations, 7 Qed., 7 Admitted.; cross-formula consistency
  relations and warning checks.
  - Unitarity.v: 23 declarations, 5 Qed., 2 Admitted.; CKM/PMNS unitarity checks and
  related candidate relations.
  - Catalog42.v: 22 declarations, 13 Qed., 0 Admitted.; representative catalogue
  aggregation file.

  The formula IDs explicitly present in the Trinity Coq layer are: C01, C02, C03, G01,
  G02, G03, G04, G06, H01, H02, H03, L01, L02, L03, N01, N03, N04, Q01, Q02, Q03,
  Q04, Q05, Q06, and Q07. Therefore the correct reading is three-tiered: the article
  restores a 42-row catalogue; the Coq layer directly mentions 24 formula IDs and
  provides 39 ID-named theorem/lemma declarations; the broader evaluator/proof
  infrastructure is larger than 42 and includes the 69-formula v0.9 interface claim, but
  not yet a closed one-theorem-per-formula proof set for all 69 entries.

  8.5 Coverage Functional Analysis
  The empirical coverage functional

  rho(C, $\varepsilon$) = (1)/(|T|) sumX $\in$ T 1[dC(X) < $\varepsilon$] 32

  is evaluated over a grid of (C, $\varepsilon$) values: C $\in$ {1,…,12}, $\varepsilon$ $\in$ {10-4, 10-3, 10-2}.
  Preliminary analysis indicates:

  - For C $\leq$ 3: coverage is consistent with random expectation under all null models.
  - For C $\in$ [4, 8]: a non-trivial excess rho(C, 10-3) > rho0(C, 10-3) is observed.
  - For C > 10: rho $\rightarrow$ 1 under saturation, and comparison with null models becomes
  uninformative.

  The critical regime C $\in$ [4, 8] at tolerance $\varepsilon$ = 10-3 is the primary domain of interest. All
  excess-coverage claims are presented as provisional pending full permutation
  calibration.

  8.6 Structural Constant alpha$\varphi$ and Experimental Test
  The quantity alpha$\varphi$ := $\varphi$-3/2 approx 0.11803 is the lowest-complexity grammar
  element near alphas(mZ). Its proximity to the PDG 2024 central value \citep{pdg2024} alphas(mZ) =
  0.1180 ± 0.0009 yields

  Deltaalpha := (|alpha$\varphi$ - alphas(mZ)|)/(alphas(mZ)) approx 2.5 $\times$ 10-4, 33

  a relative deviation well within the current experimental uncertainty. Whether this
  constitutes a statistically significant match --- after Bonferroni correction over |S(3)|
  and the full dataset --- depends on the null distribution at C = 3 and $\varepsilon$ = 2.5 $\times$ 10-4. The
  quantity alpha$\varphi$ is a low-complexity algebraic invariant arising from the grammar, not a
  derivation of alphas.

